\section{Conclusion}\label{sec:conclusion}
%==============================================================================

Exact consistency in modifiable multi-location encodings yields a deterministic finite converse family. Once the observation transcript leaves a nontrivial ambiguity class, the same finite obstruction can be stated as a confusability-clique bound, a counting converse, a conditional-entropy lower bound, and decoder-output or finite-gap constraints. The same model also supports deterministic data-processing and a budgeted finite-error extension.

The threshold $C_0=1$ is the structural corollary of that converse family: one authoritative location with deterministic views is the unique zero-incoherence regime, and any higher independent rate makes ambiguity reachable. The same obstruction then reappears operationally as manual update cost, yielding $O(1)$ cost at rate $1$ and $\Omega(n)$ cost above it, with an unbounded asymptotic gap. The realizability theorem identifies causal propagation and provenance observability as the host-level conditions under which the rate-$1$ regime is actually attained.

\textbf{Limitations and future work.} The present model is finite and exact. Natural extensions include probabilistic coherence, networked encodings with communication constraints, and partial-correlation regimes that move closer to classical stochastic side-information models.

\subsection{Artifacts}\label{sec:data-availability}

The Lean 4 formalization is included as supplementary material. Sections~\ref{sec:evaluation} and~\ref{sec:empirical} record theorem-level realizability checks and case-study material.

\subsection*{Acknowledgment: AI-use Disclosure}

Generative AI tools (including Codex, Claude Code, Augment, Kilo, and OpenCode) were used throughout this manuscript, across all sections (Abstract, Introduction, theoretical development, proof sketches, applications, conclusion, and appendix) and across all stages from initial drafting to final revision. The tools were used for boilerplate generation, prose and notation refinement, \LaTeX{}/structure cleanup, translation of informal proof ideas into candidate formal artifacts (Lean/\LaTeX{}), and repeated adversarial reviewer-style critique passes to identify blind spots and clarity gaps.

The author retained full intellectual and editorial control, including problem selection, theorem statements, assumptions, novelty framing, acceptance criteria, and final inclusion/exclusion decisions. No technical claim was accepted solely from AI output. Formal claims reported as machine-verified were admitted only after Lean verification (no \texttt{sorry} in cited modules) and direct author review; Lean was used as an integrity gate for responsible AI-assisted research. The author is solely responsible for all statements, citations, and conclusions.
